\documentclass[11pt,a4paper]{report}
\usepackage[utf8]{inputenc}
\usepackage[T1]{fontenc}
\usepackage{lmodern}
\usepackage{geometry}
\usepackage{hyperref}
\usepackage{amsmath}
\usepackage{booktabs}
\geometry{margin=2.5cm}

\title{SanaVerkko: Codebase and Algorithms\\\large with Historical Background on Gematria, Numerology, POS, and Neural Language Models}
\author{Project Documentation}
\date{\today}

\begin{document}
\maketitle
\tableofcontents

\chapter{Introduction}
SanaVerkko is an experimental system that combines symbolic text transformation, numerological feature engineering, part-of-speech (POS) constraints, and real-time audio synthesis. The project couples a simulation/control loop with a reference database of words and a ranking model that can optionally be guided by a learned long-term memory (LTM) language model.

At a high level, each word in an active sentence is represented as a node with an activation value. When activation crosses configurable thresholds, the word may be replaced by another reference word according to gematria-based matching, optional POS constraints, optional exploratory jump logic, and optional LTM-based re-ranking.

\chapter{Codebase Architecture and Main Components}
\section{Primary modules}
The codebase is centered around three primary modules:
\begin{itemize}
  \item \textbf{sanaVerkkoCore.py}: main controller, UI logic (wxPython), simulation loop, replacement logic, indexing, and integration with LTM.
  \item \textbf{sanasyna.py}: audio engine and waveform rendering; melody generation, ADSR shaping, and multi-voice counterpoint mixing.
  \item \textbf{sv\_ltm.py}: trainable next-word model (Char-CNN + MLP + softmax), serialization to \texttt{.svltm}, and inference API used by the core app.
\end{itemize}

\section{Execution model}
The runtime model is event-driven:
\begin{enumerate}
  \item UI parameters are set in wxPython controls.
  \item A timer calls \texttt{simulationStep()} periodically.
  \item Word activations update and trigger replacement checks.
  \item Updated sentence state is rendered and written to output.
  \item Audio state is regenerated when sentence/activation signature changes.
\end{enumerate}

\chapter{Core Algorithms in SanaVerkko}
\section{Gematria and numerological features}
For a word $w = c_1 c_2 \dots c_n$, gematria is computed as
\[
G(w)=\sum_{i=1}^{n} v(c_i),
\]
where $v(c_i)$ is the configured numeric value of character $c_i$.

Two derived features are used:
\begin{itemize}
  \item \textbf{Reduction}: one-pass digit sum (e.g., $137 \rightarrow 11$).
  \item \textbf{Digital root}: repeated reduction to one digit (e.g., $137 \rightarrow 11 \rightarrow 2$).
\end{itemize}

\section{Indexed candidate retrieval}
To avoid linear scans over large databases, the reference set is indexed by:
\begin{itemize}
  \item gematria,
  \item reduction,
  \item digital root,
  \item and POS-crossed versions of those indices when POS matching is enabled.
\end{itemize}

\section{Strict vs fluid gematria behavior}
The current logic supports two modes:
\begin{itemize}
  \item \textbf{Fluid gematria OFF}: strict word-level gematria is enforced; candidates are filtered to exact gematria matches only.
  \item \textbf{Fluid gematria ON}: reduction/root expansions are allowed, and jump exploration may broaden the search radius.
\end{itemize}

\section{POS matching and fluid POS}
When POS matching is enabled, retrieval is POS-conditioned.
\begin{itemize}
  \item \textbf{Fluid POS OFF}: seed POS tags are frozen from the original sentence word slots.
  \item \textbf{Fluid POS ON}: POS is recomputed dynamically from current word forms.
\end{itemize}

To avoid stagnation in narrow POS buckets, logic includes fallback broadening if no non-self alternatives exist.

\section{Candidate scoring and selection}
For a source word $s$ and candidate $c$, the base structural score is composed from gematria and numerological distances:
\[
\Delta_G = \frac{|G(c)-G(s)|}{1000},\quad
\Delta_R = \frac{|R(c)-R(s)|}{9},\quad
\Delta_D = \frac{|D(c)-D(s)|}{9},
\]
\[
\text{base}(s,c)=\Delta_G+0.35\Delta_R+0.35\Delta_D.
\]

If LTM probabilities are available, a blended score is used:
\[
\text{blend}(s,c)=(1-\lambda)\,\text{base}(s,c)+\lambda\,(1-p_{\text{ltm}}(c|\text{context})),
\]
where $\lambda \in [0,1]$ is \texttt{ltm\_weight} and $p_{\text{ltm}}$ can be reduced by a common-word penalty.

Selection combines ranked deterministic choice with controlled stochasticity (top-$k$ exploration and optional jump mode).

\section{Sentence iteration and convergence guard}
During each logic phase, words with activation magnitude above threshold are eligible for replacement. Iteration is bounded by:
\begin{itemize}
  \item maximum iteration count,
  \item and cycle detection via previously seen sentence states.
\end{itemize}
This prevents infinite loops and gives practical convergence behavior.

\section{Audio synthesis algorithm}
The sentence is converted to note sequences and rendered as looped audio:
\begin{itemize}
  \item waveform choice (sine, triangle, square, sawtooth, noise-driven modes),
  \item frequency mapping mode (original, Pythagorean, equal-tempered modal variants),
  \item polyphonic voice count/spread,
  \item ADSR envelope,
  \item minimum note duration and duration scaling.
\end{itemize}

\chapter{Gematria and Numerology in Historical Perspective}
\section{From ancient number-symbolism to formal systems}
Number symbolism appears across many ancient cultures, including Egypt and Mesopotamia, and later in Greek philosophical traditions where number harmony was treated as a structural principle of cosmos and proportion. In late antiquity and medieval religious-intellectual contexts, alphabetic-number correspondences became increasingly formalized in scriptural interpretation traditions.

\section{Hebrew Kabbalah and numerological interpretation}
In Hebrew kabbalistic practice, letter-number correspondences support interpretive techniques such as gematria, where numerical values of words are compared to identify conceptual parallels. Within this worldview, language is not only semantic but also structural-symbolic.

\section{Hermetic Qabalah (based on \texttt{chicken\_qabalah.txt})}
The source text \texttt{chicken\_qabalah.txt} frames ``Chicken Qabalah'' as a practical adaptation of Hebrew Qabalah for Western Hermetic practitioners. Several themes are emphasized:
\begin{itemize}
  \item The distinction \textit{Kabbalah/Cabala/Qabalah} is presented as partly traditional/orthodox vs practical/esoteric usage.
  \item The intended audience includes Western Hermetic domains (Tarot, ceremonial magick, Rosicrucianism, alchemy, esoteric Freemasonry, etc.).
  \item Practice is centered on letter meanings and numbers, the Tree of Life, and related meditative/symbolic exercises.
  \item The \textit{Sepher Yetzirah} is highlighted as foundational, particularly for the concepts of the ten \textit{Sephiroth} and the symbolic role of alphabetic creation.
\end{itemize}

This Hermetic framing is directly relevant to SanaVerkko because the software operationalizes a pragmatic, transformation-oriented form of gematria rather than a historical-philological reconstruction.

\chapter{Part-of-Speech (POS) Tagging}
\section{POS in NLP}
POS tagging assigns grammatical categories (noun, verb, adjective, adverb, etc.) to word tokens. In modern NLP, POS can be estimated by rule-based systems, statistical sequence models, or neural contextual models. POS constraints are useful when one wants substitutions that preserve grammatical role.

\section{POS implementation in SanaVerkko}
SanaVerkko uses a hybrid strategy:
\begin{itemize}
  \item If NLTK tagger resources are available, it uses NLTK POS tagging.
  \item Otherwise, it falls back to a suffix-based heuristic:
  \begin{itemize}
    \item \texttt{-ly} $\rightarrow$ adverb,
    \item \texttt{-ing}/\texttt{-ed} $\rightarrow$ verb,
    \item adjective-like suffixes (\texttt{-ous}, \texttt{-ful}, \texttt{-ive}, \dots),
    \item noun-like suffixes (\texttt{-ness}, \texttt{-tion}, \texttt{-ment}, \dots),
    \item default noun.
  \end{itemize}
\end{itemize}

The system caches POS tags and supports both frozen-seed POS and fluid POS mutation modes.

\chapter{Neural Language Models: General and SanaVerkko-Specific}
\section{General overview}
A neural language model estimates token probabilities conditioned on context:
\[
P(w_t\mid w_{t-k},\dots,w_{t-1}).
\]
Classic families include feed-forward $n$-gram-style MLPs, RNN/LSTM models, and Transformer-based architectures. Key training objective is usually cross-entropy minimization over next-token prediction.

\section{SanaVerkko LTM model architecture}
The model in \texttt{sv\_ltm.py} is a compact next-word predictor:
\begin{enumerate}
  \item Character-level encoder (Char-CNN) builds a fixed feature vector for each word.
  \item Context window (default size 3) concatenates recent word feature vectors.
  \item MLP head with two hidden layers projects to vocabulary logits.
  \item Softmax yields next-word probabilities.
\end{enumerate}

\section{Training objective}
The training objective is cross-entropy with L2 regularization:
\[
\mathcal{L} = -\frac{1}{B}\sum_{i=1}^{B}\log p(y_i\mid x_i)
+ \frac{\lambda}{2}\left(\|W_1\|_2^2 + \|W_2\|_2^2 + \|W_3\|_2^2\right).
\]

\section{Inference-time integration in SanaVerkko}
The core app queries LTM probabilities for candidate words and blends them into ranking with a tunable weight. This means LTM does not replace gematria logic; it re-weights candidate preference.

\section{GPU acceleration path}
Phase-1 acceleration is implemented with optional PyTorch backend:
\begin{itemize}
  \item \texttt{device=auto|cpu|mps|cuda} for training,
  \item Metal (MPS) support on Apple Silicon when available,
  \item CPU fallback when PyTorch or requested backend is unavailable.
\end{itemize}
At runtime, backend status can be displayed in the app LTM status label.

\chapter{Discussion and Future Directions}
\section{Strengths}
\begin{itemize}
  \item Clear symbolic core (gematria + numerological transforms).
  \item Practical controllability via strict/fluid toggles and exploration parameters.
  \item Efficient indexed retrieval for large reference databases.
  \item Optional neural prior (LTM) without losing interpretable symbolic behavior.
\end{itemize}

\section{Limitations}
\begin{itemize}
  \item Linguistic quality depends on reference corpus and POS coverage.
  \item Heuristic POS fallback is intentionally simple.
  \item LTM is a compact model and not as expressive as large transformer LMs.
  \item Numerological mappings are symbolic/interpretive, not scientific semantics.
\end{itemize}

\section{Potential extensions}
\begin{itemize}
  \item Richer POS/morphology constraints.
  \item Subword or transformer-based optional LTM.
  \item Better observability dashboards for candidate pools and fallback reasons.
  \item Preset metadata and experiment reproducibility tooling.
\end{itemize}

\chapter*{Source Notes}
\addcontentsline{toc}{chapter}{Source Notes}
\begin{itemize}
  \item Implementation basis: \texttt{sanaVerkkoCore.py}, \texttt{sv\_ltm.py}, \texttt{sanasyna.py} in this repository.
  \item Hermetic Qabalah framing basis: \texttt{chicken\_qabalah.txt} sections discussing practical Western Hermetic usage, Tree of Life focus, and \textit{Sepher Yetzirah} context.
\end{itemize}

\end{document}
